\documentclass[11pt]{article}

\usepackage[margin=1in]{geometry}
\usepackage[T1]{fontenc}
\usepackage[utf8]{inputenc}
\usepackage{lmodern}
\usepackage{microtype}
\usepackage{amsmath,amssymb,amsfonts,bm,mathtools}
\usepackage{booktabs}
\usepackage{siunitx}
\usepackage{natbib}
\usepackage{xcolor}
\usepackage{hyperref}
\usepackage{lineno}
\usepackage{enumitem}

\hypersetup{colorlinks=true,citecolor=blue,linkcolor=blue,urlcolor=blue}
\sisetup{detect-all}
\linenumbers

\newcommand{\dd}{\mathrm{d}}
\newcommand{\Dt}{\frac{\mathrm{D}}{\mathrm{D}t}}
\newcommand{\tr}{\operatorname{tr}}
\newcommand{\sgn}{\operatorname{sgn}}
\newcommand{\Id}{\bm{I}}
\newcommand{\vel}{\bm{u}}
\newcommand{\pos}{\bm{r}}
\newcommand{\Tvec}{\bm{T}}
\newcommand{\Kvec}{\bm{K}}
\newcommand{\pvec}{\bm{p}}
\newcommand{\nvec}{\bm{n}}
\newcommand{\Lten}{\bm{L}}
\newcommand{\Dten}{\bm{D}}
\newcommand{\Ften}{\bm{F}}
\newcommand{\Jten}{\bm{J}}
\newcommand{\Hten}{\bm{H}}
\newcommand{\eps}{\varepsilon}
\newcommand{\todo}[1]{\textcolor{red}{[TODO: #1]}}

\title{From Local Wrinkle Criteria to Finite-Width Fold Growth in Glacier Stratigraphy}
\author{Brad P. Lipovsky\\
\small Department of Earth and Space Sciences, University of Washington, Seattle, Washington, USA\\
\small Correspondence: \texttt{bradlip@uw.edu}}
\date{}

\begin{document}
\maketitle

\begin{abstract}
Stratigraphic folding limits the fidelity of ice-core climate records and records otherwise inaccessible changes in glacier flow. This tension is especially acute in the Allan Hills Blue Ice Area, East Antarctica, where shallow cores contain ice and air extending from the Pleistocene into the Pliocene and Miocene, yet exhibit strong thinning, age discontinuities, and folding. Existing theory for passive folding in ice sheets was established by Waddington, Bolzan and Alley (WBA), who derived a local criterion comparing the rotation of an infinitesimal layer disturbance with that of a steady reference isochrone. Subsequent work incorporated anisotropic rheology and finite strain along particle paths, but the fundamental diagnostic remains an orientation criterion for a zero-length tangent element. Here we formulate stratigraphic folding as the finite deformation of an entire material curve. We first show that the WBA criterion is the projective dynamics of the material-line equation and can be written as an eigendirection problem for the velocity-gradient tensor in a frame rotating with the reference isochrone. We then derive an exact evolution equation for the signed curvature $\kappa$ of a two-dimensional material layer. In incompressible flow,
\[
\Dt\kappa=-3(\pvec\!\cdot\!\Dten\pvec)\kappa
+\nvec\!\cdot\!\left[(\nabla\nabla\vel):(\pvec\otimes\pvec)\right],
\]
where $\pvec$ and $\nvec$ are the unit tangent and normal and $\Dten$ is the strain-rate tensor. Tangential compression amplifies existing curvature, whereas spatial variation of the velocity gradient creates curvature from an initially straight layer. Fold sharpening is diagnosed by $\kappa\Dt\kappa>0$, while overturning occurs when the horizontal projection of the material tangent changes sign. The framework distinguishes local rotation, finite-width curvature growth, passive finite-time amplification, and active rheological instability. It also provides a natural basis for testing whether traveling basal-traction or surface-ablation anomalies can generate the folded and discontinuous stratigraphy observed in blue-ice archives.
\end{abstract}

\section{Introduction}

Ice-sheet stratigraphy is both a climate archive and a deformation record. In the simplest setting, annual and climatic layers are deposited at the surface, advected into the ice, and progressively thinned by vertical compression. This ordered picture underpins the interpretation of deep ice cores and radar-derived age fields. Yet folded, overturned, and locally repeated layers are widespread in glacier ice, from millimetre-scale cloudy bands in cores to kilometre-scale internal-reflection structures in radar data \citep{Hudleston2015,Bons2016,Bons2025,Zhang2024}. Such structures create two complementary scientific problems. First, they can reorder or omit climatic intervals and therefore compromise chronology. Second, they retain information about transient basal traction, divide migration, anisotropy, and other dynamical processes that are weakly expressed at the present-day surface.

The paleoclimate stakes are particularly high in Antarctic blue-ice areas. Ablation and upward ice motion bring ancient ice close to the surface, permitting access to ice substantially older than the oldest continuous deep core \citep{WhillansCassidy1983,Bintanja1999,Spaulding2012}. At the Allan Hills, shallow cores have yielded atmospheric and temperature information near \SI{1}{Ma} \citep{Higgins2015}, discontinuous ice older than \SI{2}{Ma} \citep{Yan2019}, and recently dated Pliocene and Miocene ice extending to approximately \SI{6}{Ma} \citep{Shackleton2025}. These records now support reconstructions of greenhouse gases and ocean heat content over intervals previously accessible only through indirect marine and terrestrial proxies \citep{MarksPeterson2026,Shackleton2026}. Their interpretation, however, is inseparable from ice dynamics: the Allan Hills archive is strongly thinned, discontinuous, and locally folded, with age reversals over short core intervals \citep{Higgins2015,Hudak2026,Stoll2026}.

A foundational theory for passive stratigraphic folding was developed by \citet{Waddington2001}, hereafter WBA. Their central question was local and kinematic: given a pre-existing disturbance in a material layer, does the ambient flow rotate the disturbance back toward the steady stratigraphy, or away from it toward overturning? Their answer was a critical-slope criterion governed by the ratio of shear to longitudinal deformation. The approach was extended to strongly anisotropic ice \citep{ThorsteinssonWaddington2002}, to finite deformation along particle paths \citep{JacobsonWaddington2004}, and to folds generated by unsteady divide migration \citep{JacobsonWaddington2005}. Related work has demonstrated that traveling basal-slip anomalies can produce thickness-scale folds \citep{Wolovick2014,WolovickCreyts2016} and that anisotropy and convergent flow can produce large-amplitude fold systems \citep{Bons2016,Zhang2024}.

The WBA framework is powerful because it isolates the competition between shear-driven rotation and extension-driven flattening. Its limitation is equally clear: a finite-width fold is represented only through the orientation of an infinitesimal tangent segment. The criterion does not directly contain fold amplitude, wavelength, curvature, or the spatial variation of deformation across the two limbs of a fold. These omissions do not invalidate the local theory, but they obscure the distinction between four questions:
\begin{enumerate}[leftmargin=*]
    \item Does one tangent rotate away from a reference isochrone?
    \item Does an entire finite-width wrinkle sharpen or flatten?
    \item Does the wrinkle overturn and reverse stratigraphic order along a sampling line?
    \item Does an infinitesimal perturbation actively grow through feedback on the velocity field?
\end{enumerate}

Here we develop a geometric generalization of WBA that answers the first three questions for passive stratigraphy in a prescribed flow. We formulate an isochrone as a material curve, derive exact tangent and curvature evolution equations, and identify separate conditions for curvature amplification and overturning. The local WBA result emerges as the orientation dynamics of the tangent equation. The finite-width correction enters through the spatial gradient of the velocity-gradient tensor, equivalently the velocity Hessian. We then outline a numerical program for applying the theory to traveling basal-traction and surface-mass-balance anomalies and to the complex flow geometry of the Allan Hills.

\section{Motivation from the Allan Hills old-ice archive}

\subsection{Ancient ice in an ablation setting}

Blue-ice areas are mass-loss regions embedded within the broader accumulation zone of an ice sheet. Persistent wind scour and sublimation remove surface mass, while topographic obstruction and ice-flow convergence drive old ice upward toward the surface \citep{Bintanja1999,Spaulding2012}. The Allan Hills Blue Ice Area lies adjacent to the Transantarctic Mountains and contains several exposed ice fields separated by snow-covered accumulation zones. Geodetic measurements show slow horizontal motion and a resolvable upward component in the main ice field, providing the kinematic basis for surface exposure of ancient ice \citep{Spaulding2012}. Radar mapping and flowline modeling further indicate that the duration and continuity of the archive vary sharply with position and basal geometry \citep{Kehrl2018}.

The scientific value of this archive has expanded rapidly. Horizontal and vertical cores first established the presence of climate records spanning approximately 90--250 ka \citep{Spaulding2013}. A later shallow core recovered a stratigraphically disturbed section near \SI{1}{Ma}, extending direct atmospheric measurements across the Mid-Pleistocene Transition \citep{Higgins2015}. Subsequent work recovered discontinuous ice older than \SI{2}{Ma} \citep{Yan2019}. More recent cores contain individually dated samples interpreted to extend through the Pliocene and into the Miocene, with maximum ages near \SI{6}{Ma} \citep{Shackleton2025}. These samples have enabled estimates of atmospheric greenhouse gases and mean ocean temperature over much of the past three million years \citep{MarksPeterson2026,Shackleton2026}.

\subsection{Stratigraphic complexity as both limitation and signal}

The same flow processes that expose old ice also complicate the archive. Upward flow, strong cumulative thinning, spatially varying ablation, bedrock relief, and possible changes in basal traction can compress, truncate, and fold isochrones. In the Allan Hills cores, old ice is commonly present as discontinuous snapshots rather than a single monotonic age--depth sequence \citep{Yan2019,Shackleton2025}. High-resolution continuous-flow measurements reveal abrupt age discontinuities and folding, and suggest that the sampled climate-state distribution may be biased by differential preservation and deformation \citep{Hudak2026}. Microstructural observations additionally show complex fabrics and textures in million-year-old ice, raising the possibility that anisotropic deformation both records and modifies fold growth \citep{Stoll2026}.

These observations motivate a theory that can be evaluated directly against core and radar geometry. A useful model should distinguish at least three mechanisms. First, a stationary ablation zone produces upward-bending layers that terminate at an erosional unconformity. Second, a traveling surface-ablation or basal-traction anomaly can phase-lock with ice at a particular depth and greatly increase residence time in a deforming region. Third, rheological anisotropy or layer-scale viscosity contrasts can generate an active folding instability with its own wavelength dependence. The WBA theory addresses passive rotation of an existing disturbance. The development below retains that insight while introducing the finite-width geometry needed to distinguish sharpening, overturning, and truncation.

\section{Background: local theories of passive folding}

\subsection{Material-line evolution}

Consider a two-dimensional velocity field
\[
\vel(x,z,t)=\begin{pmatrix}u(x,z,t)\\w(x,z,t)\end{pmatrix},
\qquad
\Lten=\nabla\vel=
\begin{pmatrix}
 u_x & u_z\\
 w_x & w_z
\end{pmatrix}.
\]
Let $\delta\pos$ be an infinitesimal vector joining two neighboring material particles. A first-order Taylor expansion of the velocity gives the standard material-line equation
\begin{equation}
\Dt\delta\pos=\Lten\,\delta\pos.
\label{eq:material_line}
\end{equation}
Equation~\eqref{eq:material_line} is general continuum kinematics rather than a glacier-specific constitutive result \citep{Batchelor1967,Ottino1989}. WBA used it to describe the local evolution of an isochrone segment.

Write
\[
\delta\pos=\begin{pmatrix}\delta x\\\delta z\end{pmatrix},
\qquad
m=\frac{\delta z}{\delta x}.
\]
Applying the quotient rule to Equation~\eqref{eq:material_line} yields
\begin{equation}
\Dt m=w_x+(w_z-u_x)m-u_zm^2.
\label{eq:riccati}
\end{equation}
This is a Riccati equation: the apparently nonlinear slope dynamics are the projective representation of the linear vector equation~\eqref{eq:material_line}. The divergence of $m$ at a vertical tangent is a coordinate singularity; the material vector itself remains regular.

If the segment angle is $\theta$, with unit tangent and normal
\[
\pvec=(\cos\theta,\sin\theta),
\qquad
\nvec=(-\sin\theta,\cos\theta),
\]
then its angular velocity is
\begin{equation}
\Dt\theta=\nvec\cdot\Lten\pvec.
\label{eq:theta}
\end{equation}
The antisymmetric part of $\Lten$ contributes rigid-body rotation, while the symmetric part changes line orientation relative to the principal strain directions.

\subsection{The WBA critical orientation}

Let $m_0(x,z)$ denote the slope of the steady reference isochrone passing through a spatial point. Steady does not imply that a material segment on that isochrone has zero angular velocity; a particle can rotate as it traverses a spatially curved but time-independent layer. WBA therefore compared the angular velocity of a disturbed segment with that of the local reference isochrone.

Let
\[
\omega(m)=\Dt\theta(m),
\qquad
\omega_0=\omega(m_0).
\]
The instantaneous relative angle is stationary when
\begin{equation}
\omega(m)=\omega_0.
\label{eq:equal_rotation}
\end{equation}
One solution is trivially $m=m_0$. Because the equal-rotation equation is quadratic in $m$, a second solution generally exists. WBA identified this second orientation as a critical slope separating disturbances that rotate back toward the reference layer from those that rotate farther away. In the angle convention used here, its exact value is
\begin{equation}
 m_{\mathrm{c}}=-\frac{u_x-w_z}{u_z+\omega_0}-m_0,
\label{eq:mc}
\end{equation}
provided $u_z+\omega_0\neq0$. This is WBA's Equation~(6), noting that their clockwise-positive angle satisfies $\dot\theta_{0,\mathrm{WBA}}=-\omega_0$. Defining $\Delta m=m-m_0$, the exact critical disturbance is therefore
\begin{equation}
 \Delta m_{\mathrm{c}}=-\frac{u_x-w_z}{u_z+\omega_0}-2m_0.
\label{eq:dmc_exact}
\end{equation}
Near an ice divide, WBA assumed that the steady isochrone is gently sloping, so that $\omega_0\simeq w_x$. Their approximation becomes
\begin{equation}
 \Delta m_{\mathrm{c}}\simeq-\frac{u_x-w_z}{u_z+w_x}-2m_0.
\label{eq:dmc}
\end{equation}

WBA introduced the dimensionless shear number
\begin{equation}
 S=\frac{u_z+w_x}{u_x-w_z}.
\label{eq:shear_number}
\end{equation}
For two-dimensional incompressible flow, $u_x+w_z=0$, so
\begin{equation}
 S=\frac{D_{xz}}{D_{xx}},
\end{equation}
where $\Dten=(\Lten+\Lten^{\mathsf T})/2$. Thus $S$ compares shear strain rate with longitudinal extension or compression. In WBA's near-horizontal approximation,
\begin{equation}
\Delta m_{\mathrm{c}}\simeq-S^{-1}-2m_0,
\end{equation}
and for $|m_0|\ll |S|^{-1}$ this further reduces to $\Delta m_{\mathrm{c}}\simeq-S^{-1}$. Strong shear therefore permits gentle disturbances to overturn, whereas strong extension requires an initially steep disturbance. We retain the term \emph{WBA shear number}; an unqualified eponym such as ``Waddington number'' would obscure the contributions of Bolzan and Alley and the established terminology.

\subsection{Eigenstructure of the WBA criterion}

The WBA construction has a compact coordinate-free interpretation. Introduce the $90^{\circ}$ rotation tensor
\[
\Jten=\begin{pmatrix}0&-1\\1&0\end{pmatrix},
\qquad
\Jten\pvec=\nvec.
\]
Define the velocity-gradient tensor in a frame rotating with the reference isochrone,
\begin{equation}
\bm{A}=\Lten-\omega_0\Jten.
\label{eq:Arot}
\end{equation}
Then
\[
\nvec\cdot\bm{A}\pvec=\omega(\pvec)-\omega_0.
\]
A direction has the same angular velocity as the reference isochrone precisely when $\bm{A}\pvec$ has no component normal to $\pvec$. Therefore
\begin{equation}
\bm{A}\pvec=\lambda\pvec.
\label{eq:wba_eigen}
\end{equation}
The reference and critical orientations are the two eigendirections of $\bm{A}$. Under a locally constant-gradient approximation, if their eigenvalues are $\lambda_0$ and $\lambda_{\mathrm{c}}$, a small component along the critical direction evolves relative to the reference direction as
\begin{equation}
\epsilon(t)\sim\epsilon(0)\exp\left[(\lambda_{\mathrm{c}}-\lambda_0)t\right].
\label{eq:eigengap}
\end{equation}
The eigendirections define the local orientation phase portrait, while the eigenvalue gap determines attraction toward or repulsion from the reference direction. Equation~\eqref{eq:wba_eigen} is equivalent to the two roots of the Riccati equation in the co-rotating frame.

\subsection{Extensions by Waddington and collaborators}

\citet{ThorsteinssonWaddington2002} incorporated strongly anisotropic viscosity and showed that the balance between flattening and overturning depends on fabric-modified deformation, not merely isotropic Glen flow. \citet{JacobsonWaddington2004} then replaced the strictly instantaneous criterion with deformation-gradient calculations along particle paths. If
\begin{equation}
\Dt\Ften=\Lten\Ften,
\qquad
\Ften(t_0)=\Id,
\label{eq:F}
\end{equation}
then an initial tangent $\Tvec_0$ evolves as $\Tvec(t)=\Ften(t,t_0)\Tvec_0$. This finite-time construction accounts for the complete deformation history between an upstream disturbance and an ice-core site. \citet{JacobsonWaddington2005} applied related ideas to abandoned Raymond arches created by divide migration and showed that open arches can evolve into recumbent folds downstream.

These developments resolve the finite-time evolution of a tangent vector, but they do not by themselves provide a local equation for the spatial width and curvature of an entire fold. That distinction motivates the following whole-curve formulation.

\section{Curvature model for finite-width folds}

\subsection{Isochrones as material curves}

Parameterize an isochrone by a material coordinate $a$,
\begin{equation}
\pos=\pos(a,t)=\begin{pmatrix}x(a,t)\\z(a,t)\end{pmatrix},
\qquad
\partial_t\pos(a,t)=\vel(\pos(a,t),t).
\label{eq:curve_advect}
\end{equation}
Define the tangent, second material derivative, tangent stretch, unit tangent, and unit normal by
\begin{equation}
\Tvec=\partial_a\pos,
\qquad
\Kvec=\partial_a^2\pos,
\qquad
q=|\Tvec|,
\qquad
\pvec=\frac{\Tvec}{q},
\qquad
\nvec=\Jten\pvec.
\label{eq:curve_defs}
\end{equation}
Differentiating Equation~\eqref{eq:curve_advect} with respect to $a$ gives
\begin{equation}
\Dt\Tvec=\Lten\Tvec,
\label{eq:T_evolution}
\end{equation}
which is the WBA material-line equation applied pointwise along the entire layer. Differentiating again gives
\begin{equation}
\Dt\Kvec=\Lten\Kvec+
(\nabla\Lten):(\Tvec\otimes\Tvec).
\label{eq:K_evolution}
\end{equation}
In components,
\[
\Dt K_i=u_{i,j}K_j+u_{i,jk}T_jT_k.
\]
The first term stretches and rotates existing spatial variation of the tangent. The second term is the finite-width correction: neighboring parts of the layer experience different velocity gradients.

\subsection{Exact signed-curvature evolution}

The signed curvature of the material curve is
\begin{equation}
\kappa=\frac{\det(\Tvec,\Kvec)}{|\Tvec|^3}.
\label{eq:kappa_def}
\end{equation}
The tangent stretch rate is
\begin{equation}
\epsilon_t=\pvec\cdot\Dten\pvec,
\qquad
\Dten=\frac{1}{2}(\Lten+\Lten^{\mathsf T}),
\label{eq:et}
\end{equation}
so that $\Dt q=q\epsilon_t$. Using the two-dimensional identity
\[
\det(\Lten\Tvec,\Kvec)+\det(\Tvec,\Lten\Kvec)
=(\tr\Lten)\det(\Tvec,\Kvec),
\]
Equations~\eqref{eq:K_evolution}--\eqref{eq:kappa_def} yield
\begin{equation}
\boxed{
\Dt\kappa=
\left(\tr\Lten-3\epsilon_t\right)\kappa
+\nvec\cdot\left[(\nabla\nabla\vel):(\pvec\otimes\pvec)\right].
}
\label{eq:kappa_general}
\end{equation}
For incompressible ice flow, $\tr\Lten=\nabla\cdot\vel=0$, and therefore
\begin{equation}
\boxed{
\Dt\kappa=-3(\pvec\cdot\Dten\pvec)\kappa
+\nvec\cdot\left[(\nabla\nabla\vel):(\pvec\otimes\pvec)\right].
}
\label{eq:kappa_incompressible}
\end{equation}
We denote the curvature-source term by
\begin{equation}
Q_{\kappa}=\nvec\cdot\left[(\nabla\nabla\vel):(\pvec\otimes\pvec)\right]
=\nvec\cdot(\pvec\cdot\nabla)^2\vel.
\label{eq:Qk}
\end{equation}

Equation~\eqref{eq:kappa_incompressible} separates two mechanisms. The homogeneous term $-3\epsilon_t\kappa$ amplifies or suppresses curvature already present. Tangential compression, $\epsilon_t<0$, increases $|\kappa|$; tangential extension, $\epsilon_t>0$, flattens the layer. The inhomogeneous term $Q_{\kappa}$ creates curvature where the velocity gradient varies along the tangent direction. An initially straight layer, $\kappa=0$, bends at the instantaneous rate $\Dt\kappa=Q_{\kappa}$.

The factor of three has a geometric origin. One factor arises from the evolution of the numerator in Equation~\eqref{eq:kappa_def}, while three powers of the tangent stretch appear in its denominator. In a spatially affine incompressible flow, $\nabla\Lten=0$ and $Q_{\kappa}=0$; straight lines remain straight, while existing curvature is amplified by tangential compression and reduced by extension.

\subsection{Local conditions for fold sharpening}

For $\kappa\neq0$, local curvature magnitude grows when
\begin{equation}
\Dt|\kappa|>0
\quad\Longleftrightarrow\quad
\kappa\Dt\kappa>0.
\label{eq:local_growth_def}
\end{equation}
Combining this definition with Equation~\eqref{eq:kappa_incompressible} gives
\begin{equation}
\boxed{
-3\epsilon_t\kappa^2+\kappa Q_{\kappa}>0.
}
\label{eq:local_growth}
\end{equation}
Equation~\eqref{eq:local_growth} is the finite-width local generalization of the WBA orientation criterion. It states that a bend sharpens when tangential compression and differential rotation reinforce its signed curvature. It also makes clear that no single ratio analogous to $S$ can describe all finite-width folds: the source term depends on the curvature sign, tangent orientation, and spatial variation of the velocity gradient.

A useful nondimensional diagnostic for $\kappa\neq0$ is
\begin{equation}
\mathcal{G}_{\kappa}
=\frac{1}{|\dot\epsilon_*|}\frac{1}{\kappa}\Dt\kappa
=\frac{-3\epsilon_t+Q_{\kappa}/\kappa}{|\dot\epsilon_*|},
\label{eq:Gk}
\end{equation}
where $\dot\epsilon_*$ is a chosen characteristic strain rate. Positive $\mathcal{G}_{\kappa}$ denotes instantaneous exponential amplification of signed curvature; it is undefined at $\kappa=0$, where $Q_{\kappa}$ itself is the appropriate diagnostic.

\subsection{A global finite-width measure}

Local curvature growth at one point need not imply growth of the whole fold. For a material layer segment $\mathcal{C}(t)$, define curvature energy
\begin{equation}
E_{\kappa}(t)=\frac{1}{2}\int_{\mathcal{C}(t)}\kappa^2\,\dd s.
\label{eq:Ek}
\end{equation}
Because a material arclength element evolves as $\Dt(\dd s)=\epsilon_t\dd s$, differentiation gives
\begin{equation}
\boxed{
\frac{\dd E_{\kappa}}{\dd t}
=\int_{\mathcal{C}(t)}
\left[-\frac{5}{2}\epsilon_t\kappa^2+\kappa Q_{\kappa}\right]\dd s.
}
\label{eq:Ek_growth}
\end{equation}
Thus a finite fold grows in this norm when the right-hand side of Equation~\eqref{eq:Ek_growth} is positive. Other norms may be preferable for particular observations: maximum curvature for core-scale sharp hinges, total variation of tangent angle for fold trains, or vertical amplitude for radar-scale structures. The important point is that a finite-width criterion must integrate the spatially varying geometry of the whole layer.

\subsection{Overturning and stratigraphic order}

Curvature growth is not equivalent to overturning. A fold overturns relative to a vertical core or a horizontal flow coordinate when the horizontal projection of the tangent loses monotonicity. Define
\begin{equation}
J_x(a,t)=\partial_a x=\bm{e}_x\cdot\Tvec.
\label{eq:Jx}
\end{equation}
If the layer is initially ordered with $J_x>0$, the first vertical tangent occurs when
\begin{equation}
\boxed{
\min_a J_x(a,t)=0.
}
\label{eq:overturn}
\end{equation}
An interval with $J_x<0$ is locally overturned. The WBA slope $m=z_a/x_a$ diverges as $J_x\rightarrow0$, but the material curve remains regular. Equation~\eqref{eq:overturn} therefore provides a nonsingular finite-amplitude overturning criterion.

Actual age reversal in a vertical core additionally depends on the intersection geometry between the folded isochrone and the core line. A recumbent fold can create multiple intersections of one isochrone with a vertical sampling line, while an ablation surface can remove one or more limbs and replace a fold signature with a truncation or unconformity. Core interpretation consequently requires both the internal material geometry and the history of the non-material surface boundary.

\subsection{Finite-time amplification and non-normality}

Equations~\eqref{eq:kappa_incompressible} and \eqref{eq:local_growth} are instantaneous. Along a trajectory, compression and curvature production may change sign. The finite-time solution is
\begin{equation}
\kappa(t)=
\exp\left[-3\int_{t_0}^{t}\epsilon_t(\tau)\,\dd\tau\right]
\left[
\kappa(t_0)+
\int_{t_0}^{t}
\exp\left(3\int_{t_0}^{\tau}\epsilon_t(s)\,\dd s\right)
Q_{\kappa}(\tau)\,\dd\tau
\right].
\label{eq:kappa_solution}
\end{equation}
This expression separates accumulated amplification of initial curvature from curvature generated along the path.

For the complete layer perturbation, let $\mathcal{U}(t,t_0)$ denote the linearized finite-time propagator mapping an initial displacement field $\bm{\xi}(a,t_0)$ to $\bm{\xi}(a,t)$. Fold amplification in a chosen norm is governed by the singular values of $\mathcal{U}$, not generally by the instantaneous eigenvalues of $\Lten$. This distinction matters in spatially varying and non-normal flows, where transient growth can be large even when all local modal growth rates are non-positive.

\subsection{Amplitude, wavelength, and the absence of passive wavelength selection in affine flow}

Consider an initial normal displacement
\begin{equation}
\bm{\xi}(a,t_0)=A_0\cos(ka)\nvec_0,
\qquad
k=\frac{2\pi}{\lambda}.
\label{eq:sinusoid}
\end{equation}
Its slope amplitude scales as $A_0k$ and its curvature amplitude as $A_0k^2$. A local orientation criterion therefore senses amplitude divided by wavelength, whereas the curvature formulation also contains the second spatial derivative needed to distinguish broad and narrow wrinkles.

If the velocity field is affine,
\begin{equation}
\vel(\pos,t)=\Lten(t)\pos+\bm{b}(t),
\end{equation}
then every material curve is transformed by the same deformation gradient $\Ften$. No intrinsic wavelength is selected: all Fourier components are acted on by the same spatially uniform linear map. Wavelength dependence in passive folding requires spatial variation in $\Lten$, finite boundaries with a response depending on $kH$, or a forcing region of finite width. A preferred wavelength in an actively growing fold additionally requires a rheological or dynamical feedback, such as anisotropy, viscosity contrast, damage, or another intrinsic length scale \citep{Bons2025,Zhang2024}.

\section{Numerical formulation}

The theory can be tested with a hierarchy of prescribed-flow and dynamically coupled calculations. In all cases, material stratigraphy may be represented either by Lagrangian marker curves satisfying Equation~\eqref{eq:curve_advect} or by a material scalar $\tau(\bm{x},t)$ satisfying
\begin{equation}
\partial_t\tau+\vel\cdot\nabla\tau=0,
\label{eq:age_advection}
\end{equation}
whose level sets are isochrones. Lagrangian curves provide direct access to $\Tvec$, $\kappa$, and $J_x$; an Eulerian scalar is more robust after strong overturning and for three-dimensional extension.

For dynamically consistent experiments, the velocity is obtained from incompressible Stokes flow,
\begin{align}
\nabla\cdot\vel&=0,\\
-\nabla p+\nabla\cdot\bm{\tau}+\rho\bm{g}&=\bm{0},
\end{align}
with isotropic Glen rheology or an anisotropic constitutive relation as appropriate. At the upper surface $z=s(x,t)$,
\begin{equation}
\partial_t s+u_s\partial_x s-w_s=\dot a(x,t),
\label{eq:surfacekin}
\end{equation}
where $\dot a$ is surface mass balance. Basal boundary conditions may prescribe no penetration and a spatially or temporally varying traction or sliding law.

Four benchmark classes are proposed:
\begin{enumerate}[leftmargin=*]
    \item \textbf{Affine-flow verification.} Compare numerical tangent and curvature evolution against Equations~\eqref{eq:F} and \eqref{eq:kappa_solution} with $Q_{\kappa}=0$.
    \item \textbf{Spatially varying prescribed flow.} Construct analytic velocity fields with known $\nabla\Lten$ to verify curvature production from an initially straight layer.
    \item \textbf{Traveling basal-traction and ablation anomalies.} Impose perturbations proportional to $f(x-ct)$ and quantify fold amplification as functions of anomaly speed $c$, width $L$, ice thickness $H$, and depth-dependent ice velocity $u(z)$.
    \item \textbf{Allan Hills geometry.} Use measured surface velocity, mass balance, bed topography, radar layers, and dated core horizons to evaluate whether observed folds and age reversals can be generated by passive deformation alone.
\end{enumerate}

For a traveling forcing, the principal nondimensional controls are expected to include
\begin{equation}
 kH,\qquad \frac{c}{U_s},\qquad \frac{c}{U(z)},\qquad
 \frac{A L}{H|U-c|},\qquad
 \frac{u_b}{U_s},
\label{eq:dimensionless}
\end{equation}
where $A$ is an ablation-rate scale, $U_s$ is surface speed, and $u_b$ is basal speed. The factor $AL/[H|U-c|]$ measures the potential vertical displacement accumulated during residence beneath a traveling surface anomaly. Model output should report the WBA orientation tendency, $\kappa$, $Q_{\kappa}$, $E_{\kappa}$, and the first time at which $J_x=0$.

\section{Numerical modeling results}

% Intentionally left blank for numerical results, figures, and quantitative tests.

\section{Discussion}

\subsection{From tangent rotation to finite-width folding}

The WBA criterion and the curvature model address different but nested questions. WBA diagnoses the relative rotation of one infinitesimal tangent at one position. The curvature model diagnoses variation of tangent orientation along a finite layer. In this hierarchy,
\begin{equation}
\Lten\quad\longrightarrow\quad\text{rotation and stretching of one tangent},
\end{equation}
whereas
\begin{equation}
\nabla\Lten\quad\longrightarrow\quad\text{differential rotation of neighboring tangents and curvature production}.
\end{equation}
The WBA critical orientation is therefore not superseded; it is recovered as the local projective dynamics of Equation~\eqref{eq:T_evolution}. The new ingredient is an explicit account of the fold's spatial width.

This distinction resolves an apparent conceptual problem in local wrinkle analyses. A finite-amplitude bump does not possess one slope. It consists of a continuum of tangent segments $m(a,t)$, each evolving under the local velocity gradient. WBA may be applied pointwise to these segments, but the resulting collection becomes a theory of a finite fold only when the segments are coupled through their spatial arrangement. Equation~\eqref{eq:kappa_incompressible} performs that coupling at the level of curvature.

\subsection{Passive amplification versus active instability}

The present theory assumes a prescribed velocity field. It determines whether an existing disturbance is rotated, sharpened, or overturned by the flow. It does not establish that an arbitrarily small disturbance grows through feedback on stress or viscosity. An active stability theory must perturb the momentum and constitutive equations as well as the stratigraphy. If $\tau$ is a material layer label and primes denote perturbations,
\begin{equation}
\partial_t\tau'+\vel_0\cdot\nabla\tau'
+\vel'\cdot\nabla\tau_0=0.
\label{eq:active_layer}
\end{equation}
The perturbation velocity $\vel'$ must be obtained from the linearized Stokes and rheological equations. Fourier modes proportional to $\exp(ikx+\sigma t)$ are actively unstable when $\Re\sigma(k)>0$. This coupled problem can select a wavelength, unlike passive deformation in an affine flow.

Strong ice anisotropy provides one plausible feedback. Observed fold spectra and three-dimensional simulations indicate that crystal fabric and power-law rheology can promote fold formation over a broad range of scales \citep{Bons2016,Bons2025,Zhang2024}. The curvature equation remains useful in such models as a geometric diagnostic, but $\vel$ can no longer be regarded as independent of the folded stratigraphy and fabric.

\subsection{Traveling basal-traction anomalies}

Traveling basal-slip or sticky patches provide a clear example of strong passive amplification. Stationary changes in basal slip redistribute horizontal flux with depth and produce vertical motion through incompressibility. When the traction anomaly travels at speed $c$, ice for which $u(z)\simeq c$ remains near the anomaly for an anomalously long time. Thermomechanical models show that such patches can uplift basal ice through a substantial fraction of the ice thickness and produce long-lived flux perturbations \citep{Wolovick2014}. Kinematic models show that the moving-frame flow can contain closed streamlines and material trapping \citep{WolovickCreyts2016}.

In the present language, the anomaly produces both large integrated tangential strain and a strong curvature source $Q_{\kappa}$ near its leading and trailing transitions. The WBA criterion predicts whether local tangents rotate toward overturning; Equation~\eqref{eq:kappa_solution} predicts how the finite-width bend accumulates along a trajectory; and Equation~\eqref{eq:overturn} identifies actual reversal of the horizontal projection. This decomposition may clarify why some basal anomalies produce smooth layer drape while others create recumbent folds.

\subsection{Traveling ablation zones in an accumulation region}

A traveling surface-ablation zone is related but not identical. The surface is not a material streamline when $\dot a\neq0$, and material trajectories can enter the ice through accumulation and terminate at the surface through ablation. The generic response is therefore upward bending and truncation of layers, followed by burial of an erosional unconformity when accumulation resumes. Nevertheless, speed matching can create large deformation. A patch of width $L$ moving at speed $c$ has a nominal residence time $L/|u-c|$, so modest vertical velocities can produce thickness-scale displacements when $c$ approaches the local ice speed.

Whether the result is an intact fold or a truncated unconformity depends on competition among curvature growth, upward advection, and removal at the surface. The curvature criterion alone is insufficient because the material curve may cease to exist after intersecting the ablation boundary. Numerical experiments should therefore track both internal geometry and first-passage time to the surface. A useful observational distinction is that passive folding preserves continuity of a material horizon through the fold, whereas ablation creates terminations and age gaps subsequently draped by younger accumulation.

\subsection{Implications for Allan Hills stratigraphy}

The Allan Hills provide a natural test of the theory because the system combines nearly every relevant ingredient: slow flow, strong and spatially variable ablation, upward motion, rugged basal and lateral topography, long residence times, extreme cumulative thinning, and complex crystal fabric. The old-ice archive need not be explained by one universal folding mechanism. Different cores and depth intervals may sample distinct regimes:
\begin{enumerate}[leftmargin=*]
    \item passive upwarping and surface truncation associated with stationary blue-ice ablation;
    \item finite-time fold amplification as ice traverses spatially varying strain fields around bedrock obstacles;
    \item transient deformation associated with migrating ablation boundaries or changing basal traction;
    \item active anisotropic folding in highly strained old ice.
\end{enumerate}

The curvature-source term provides a direct bridge between observed velocity fields and layer geometry. Where a sufficiently resolved model of $\vel$ is available, maps of $Q_{\kappa}$ predict where straight or gently curved isochrones should begin to bend. Integrated tangential compression predicts where those bends should sharpen. Radar-constrained layers then test the forward prediction, while core age reversals and continuous-flow measurements constrain whether the modeled curves actually overturn or are truncated. This approach could help determine whether age discontinuities are local structural features, repeated limbs of folds, erosional omissions, or boundaries between dynamically distinct ice packages.

\subsection{Observational tests}

The framework suggests several measurable diagnostics. First, radar-derived curvature can be compared with modeled $Q_{\kappa}$ and accumulated $\epsilon_t$ along flow paths. Second, the sign of $J_x$ can be inferred where three-dimensional radar mapping resolves overturned limbs. Third, core-scale fold hinges and crystal fabrics can test whether passive curvature amplification or anisotropic active folding is more plausible. Fourth, repeated absolute ages on closely spaced samples can distinguish continuous recumbent folds from discontinuities produced by surface or basal removal. Finally, dense surface-velocity and strain measurements can identify present-day regions in which the WBA and curvature criteria disagree, providing targeted tests of the finite-width correction.

\subsection{Limitations and extension to three dimensions}

The derivation is two-dimensional and assumes a smooth material curve in a differentiable velocity field. Real fold systems are three-dimensional material surfaces. Their geometry is described by a shape operator or second fundamental form with two principal curvatures, and their evolution involves the surface velocity gradient and its spatial derivatives. Cross-flow permits material bypass that is impossible in a two-dimensional flowline and may weaken moving-frame trapping. Extending Equation~\eqref{eq:kappa_incompressible} to the evolution of the second fundamental form is therefore a priority.

Additional limitations include neglected fracture, recrystallization, diffusion of chemical layers, and rheological feedback. Near an ablating surface, layers are not conserved indefinitely. Near the bed, basal melting, freeze-on, and debris entrainment can create or destroy stratigraphic surfaces. These processes must be represented explicitly before the theory is applied to absolute age continuity.

\section{Conclusions}

WBA established a fundamental local criterion for passive stratigraphic folding by comparing the rotation of a disturbed layer segment with that of a steady reference isochrone. Their result can be expressed as an eigendirection problem for the velocity-gradient tensor in a frame rotating with the reference layer. This interpretation clarifies that the WBA shear number governs local orientation dynamics, not the complete geometry of a finite fold.

Treating an isochrone as a material curve yields an exact finite-width extension. In incompressible two-dimensional flow, curvature evolves according to Equation~\eqref{eq:kappa_incompressible}. Existing curvature is amplified by tangential compression and suppressed by extension; curvature is created by spatial variation of the velocity gradient along the layer. The local sharpening condition is Equation~\eqref{eq:local_growth}, a global fold-growth condition follows from Equation~\eqref{eq:Ek_growth}, and actual overturning occurs when the horizontal tangent projection first vanishes. These diagnostics separate local rotation, finite-width sharpening, overturning, and active instability.

The framework is particularly relevant to the Allan Hills, where the scientific value of multi-million-year-old ice is limited by folding, thinning, truncation, and age reversal. Coupling the geometric diagnostics developed here with radar, absolute dating, core microstructure, and full-Stokes modeling offers a route toward determining which parts of the archive preserve continuous climate information and which record transient glacier dynamics.

\section*{Data and code availability}
No new data are presented in this theoretical manuscript. Numerical implementation and model output will be archived in a public repository upon completion. \todo{Insert repository and DOI.}

\section*{Author contributions}
B.P.L. developed the conceptual framework and derivations and wrote the manuscript. \todo{Revise after coauthor contributions are established.}

\section*{Competing interests}
The author declares no competing interests.

\section*{Acknowledgements}
\todo{Add acknowledgements to the Allan Hills field, ice-core, radar, and modeling teams and relevant funding programs.}

\appendix
\section{Derivation of the curvature equation}

Starting from
\[
\kappa=\frac{N}{q^3},
\qquad
N=\det(\Tvec,\Kvec),
\qquad
q=|\Tvec|,
\]
we have
\begin{equation}
\Dt q=\frac{\Tvec\cdot\Lten\Tvec}{q}=q\,\pvec\cdot\Dten\pvec=q\epsilon_t.
\end{equation}
Using Equations~\eqref{eq:T_evolution} and \eqref{eq:K_evolution},
\begin{align}
\Dt N
&=\det(\Lten\Tvec,\Kvec)
+\det(\Tvec,\Lten\Kvec)
+\det\left(\Tvec,(\nabla\Lten):(\Tvec\otimes\Tvec)\right)\\
&=(\tr\Lten)N
+q^3\nvec\cdot\left[(\nabla\nabla\vel):(\pvec\otimes\pvec)\right].
\end{align}
Therefore
\begin{align}
\Dt\kappa
&=q^{-3}\Dt N-3Nq^{-4}\Dt q\\
&=(\tr\Lten-3\epsilon_t)\kappa
+\nvec\cdot\left[(\nabla\nabla\vel):(\pvec\otimes\pvec)\right],
\end{align}
which reduces to Equation~\eqref{eq:kappa_incompressible} when $\nabla\cdot\vel=0$.

\section{Relation between slope and whole-curve descriptions}

For a layer that can locally be written as $z=h(x,t)$,
\[
\Tvec\propto(1,m),
\qquad
m=\partial_x h.
\]
The signed curvature is
\begin{equation}
\kappa=\frac{\partial_x m}{(1+m^2)^{3/2}}.
\end{equation}
Thus WBA evolves the local orientation variable $m$, whereas the curvature model evolves its arclength derivative. A finite-width disturbance requires the field $m(x,t)$, not one scalar slope. The overturning condition $x_a=0$ corresponds to $|m|\rightarrow\infty$ in graph coordinates.

\bibliographystyle{plainnat}
\bibliography{references}

\end{document}
